% Options for packages loaded elsewhere
\PassOptionsToPackage{unicode}{hyperref}
\PassOptionsToPackage{hyphens}{url}
\documentclass[
]{article}
\usepackage{xcolor}
\usepackage[margin=1in]{geometry}
\usepackage{amsmath,amssymb}
\setcounter{secnumdepth}{-\maxdimen} % remove section numbering
\usepackage{iftex}
\ifPDFTeX
  \usepackage[T1]{fontenc}
  \usepackage[utf8]{inputenc}
  \usepackage{textcomp} % provide euro and other symbols
\else % if luatex or xetex
  \usepackage{unicode-math} % this also loads fontspec
  \defaultfontfeatures{Scale=MatchLowercase}
  \defaultfontfeatures[\rmfamily]{Ligatures=TeX,Scale=1}
\fi
\usepackage{lmodern}
\ifPDFTeX\else
  % xetex/luatex font selection
\fi
% Use upquote if available, for straight quotes in verbatim environments
\IfFileExists{upquote.sty}{\usepackage{upquote}}{}
\IfFileExists{microtype.sty}{% use microtype if available
  \usepackage[]{microtype}
  \UseMicrotypeSet[protrusion]{basicmath} % disable protrusion for tt fonts
}{}
\makeatletter
\@ifundefined{KOMAClassName}{% if non-KOMA class
  \IfFileExists{parskip.sty}{%
    \usepackage{parskip}
  }{% else
    \setlength{\parindent}{0pt}
    \setlength{\parskip}{6pt plus 2pt minus 1pt}}
}{% if KOMA class
  \KOMAoptions{parskip=half}}
\makeatother
\usepackage{graphicx}
\makeatletter
\newsavebox\pandoc@box
\newcommand*\pandocbounded[1]{% scales image to fit in text height/width
  \sbox\pandoc@box{#1}%
  \Gscale@div\@tempa{\textheight}{\dimexpr\ht\pandoc@box+\dp\pandoc@box\relax}%
  \Gscale@div\@tempb{\linewidth}{\wd\pandoc@box}%
  \ifdim\@tempb\p@<\@tempa\p@\let\@tempa\@tempb\fi% select the smaller of both
  \ifdim\@tempa\p@<\p@\scalebox{\@tempa}{\usebox\pandoc@box}%
  \else\usebox{\pandoc@box}%
  \fi%
}
% Set default figure placement to htbp
\def\fps@figure{htbp}
\makeatother
\setlength{\emergencystretch}{3em} % prevent overfull lines
\providecommand{\tightlist}{%
  \setlength{\itemsep}{0pt}\setlength{\parskip}{0pt}}
\usepackage[spanish]{babel}
\usepackage{amsmath}
\usepackage{fontspec}
\usepackage{unicode-math}
\usepackage{graphicx}
\usepackage{adjustbox}
\usepackage{tikz}
\usepackage{pgfplots}
\usepackage{booktabs}
\usetikzlibrary{3d,babel}
\usepackage{bookmark}
\IfFileExists{xurl.sty}{\usepackage{xurl}}{} % add URL line breaks if available
\urlstyle{same}
\hypersetup{
  hidelinks,
  pdfcreator={LaTeX via pandoc}}

\author{}
\date{\vspace{-2.5em}}

\begin{document}

```

\section{Question}\label{question}

En la gráfica se muestra la probabilidad de que la variable aleatoria
\(x\) tome valores en cada uno de tres intervalos en que se ha dividido
la curva.

\includegraphics[width=12cm,height=\textheight,keepaspectratio]{prob_dist_grafico.png}

¿Cuál de las siguientes tablas representa la probabilidad de que la
variable aleatoria \(x\) tome los valores en el intervalo indicado?

\subsection{Answerlist}\label{answerlist}

\begin{itemize}
\item
  \includegraphics[width=0.7\linewidth,height=\textheight,keepaspectratio]{tabla_opcion_a.png}
\item
  \includegraphics[width=0.7\linewidth,height=\textheight,keepaspectratio]{tabla_opcion_b.png}
\item
  \includegraphics[width=0.7\linewidth,height=\textheight,keepaspectratio]{tabla_opcion_c.png}
\item
  \includegraphics[width=0.7\linewidth,height=\textheight,keepaspectratio]{tabla_opcion_d.png}
\end{itemize}

\section{Solution}\label{solution}

Para encontrar la tabla correcta, debemos extraer la información
directamente del gráfico y compararla con las opciones proporcionadas.

\textbf{Paso 1: Analizar la información del gráfico.}

El gráfico está dividido en tres intervalos con sus respectivas
probabilidades:

\begin{enumerate}
\def\labelenumi{\arabic{enumi}.}
\tightlist
\item
  El primer intervalo va desde \(x=0\) hasta \(x=6\). La probabilidad en
  esta sección es \textbf{0,25}.
\item
  El segundo intervalo (central) va desde \(x=6\) hasta \(x=9\). La
  probabilidad en esta sección es \textbf{0,50}.
\item
  El tercer intervalo va desde \(x=9\) hasta \(x=14\). La probabilidad
  en esta sección es \textbf{0,25}.
\end{enumerate}

\textbf{Paso 2: Construir la tabla esperada.}

Con la información anterior, la tabla correcta debe tener la siguiente
estructura:

Intervalo

Probabilidad

\(0 \le x \le 6\)

0,25

\(6 < x \le 9\)

0,50

\(9 < x \le 14\)

0,25

\textbf{Paso 3: Comparar con las opciones.}

Al revisar las tablas de las opciones, la única que coincide
perfectamente con los datos extraídos del gráfico es la tabla correcta.
Las otras tablas presentan errores comunes:

\begin{itemize}
\tightlist
\item
  \textbf{Error de probabilidad acumulada:} Una de las tablas muestra
  cómo la probabilidad se va sumando (0,25, 0,75, 1,00), lo cual es una
  interpretación incorrecta del gráfico.
\item
  \textbf{Error de definición de intervalo:} Otra tabla usa intervalos
  acumulativos (ej: \(0 \le x \le 9\)), lo cual no corresponde a las
  secciones individuales mostradas.
\item
  \textbf{Error de asignación:} Otra tabla intercambia los valores de
  probabilidad, asignando el valor central a los laterales y viceversa.
\end{itemize}

Por lo tanto, la respuesta correcta es la que representa fielmente los
tres intervalos y sus probabilidades tal como se muestran en el gráfico.

\subsection{Answerlist}\label{answerlist-1}

\begin{itemize}
\tightlist
\item
  Correcto. Esta tabla representa exactamente los tres intervalos y sus
  probabilidades correspondientes como se muestra en el gráfico.
\item
  Incorrecto. Usa intervalos acumulativos (0 \(\le\) x \(\le\) a, 0
  \(\le\) x \(\le\) b, 0 \(\le\) x \(\le\) c) y una columna acumulada;
  no corresponde a las tres secciones individuales.
\item
  Incorrecto. Intercambia los valores asignando el central a los
  laterales.
\item
  Incorrecto. No coincide con la asignación de probabilidades del
  gráfico.
\end{itemize}

\section{Meta-information}\label{meta-information}

exname: probabilidad\_distribucion\_grafico\_tabla extype: schoice
exsolution: 1000 exshuffle: TRUE exsection:
Estadística/Probabilidad/Interpretación de gráficos

\end{document}
